\subsection{Interpreting LIFE nitrogen signatures and abundance-region probabilities}
\label{sec:life_interpretation}

The recovered $\mathrm{H_2O}$, $\mathrm{CO_2}$, $\mathrm{CH_4}$, and, at longer exposure, $\mathrm{O_3}$ provide atmospheric context for evaluating habitability and biologically relevant chemical cycles. Their detection alone does not establish habitable surface conditions or a biological origin. \citet{woganSurfaceFluxes2026} emphasize the role of photochemistry and climate in connecting atmospheric abundances to the sources sustaining them. Our retrievals constrain effective abundances and leave that source attribution conditional on the adopted atmospheric and photochemical models.

Large biological $\mathrm{N_2O}$ fluxes can occur when denitrification produces $\mathrm{N_2O}$ efficiently while its final reduction to $\mathrm{N_2}$ is limited \citep{schwietermanEvaluatingPlausibleRange2022}. Enhanced $\mathrm{N_2O}$ can therefore arise in natural biospheres as well as under agricultural forcing. Atmospheric $\mathrm{NH_3}$ responds to a different balance: dissolution, conversion to $\mathrm{NH_4^+}$, deposition, and uptake by surface reservoirs can remove it efficiently, while accumulation becomes more favorable as these sinks approach saturation \citep{huangAssessmentAmmoniaBiosignature2022a}. Joint constraints on the two molecules could help test these source--sink balances, provided that plausible natural and agricultural mechanisms are compared explicitly.

The region probabilities in \autoref{sec:life_abundance_regions} quantify support for composition thresholds. Their $K_R$ values compare complementary abundance regions using the adopted prior within each region. A probability of ExoFarm would additionally require predictive distributions for agricultural and natural atmospheres and contextual prior probabilities for those hypotheses. The A0--A3 scenarios provide controlled reference calculations, but they do not define these distributions by themselves. The larger joint $K_R$ also incorporates the strong $\mathrm{N_2O}$ constraint and should not be read as equally strong, independent support from $\mathrm{NH_3}$. The A1 threshold was selected retrospectively after examining this series; it should be fixed before assessing future realizations.

Several limitations affect the numerical interpretation. The synthetic spectra contain vertical chemical structure, while the retrieval represents each molecule by a constant abundance; a pressure-weighted reference need not equal the abundance favored by the spectral contribution functions. Each dataset uses one noise realization with diagonal uncertainties, and A1--1 h uses the second realization selected after inspecting both results. The series therefore does not establish detection completeness or the coverage of the reported intervals. Verification of the complete effective prior remains necessary before treating the nominal region factors as final evidence ratios.

A subsequent interface diagnosis also identified different spectral-bin integration conventions in the generator and likelihood used for this 1--100 h series. This discrepancy can bias recovered abundances and evidence comparisons; the reported values describe the original series and require reassessment with a common integration operator before drawing final instrumental forecasts. The low recovery of the A1 effective reference reinforces the need for these controls. The A3--1000 h analysis is being treated separately, and its accepted products should be incorporated with their corresponding matched comparisons.

